\documentclass[11pt]{article}
\usepackage[margin=1in]{geometry}
\usepackage{amsmath,amssymb,amsthm,booktabs,enumitem}
\usepackage[hidelinks]{hyperref}
\newtheorem{definition}{Definition}
\newtheorem{proposition}{Proposition}
\newtheorem{theorem}{Theorem}
\newcommand{\R}{\mathbb{R}}
\newcommand{\1}{\mathbf{1}}
\DeclareMathOperator*{\argmax}{arg\,max}

\begin{document}

{\Large\bfseries Method: rules, exceptions, and how we measure them\par}\vspace{0.8em}

\noindent
Our analysis starts from one object, the \emph{answer table} of a task, and from it derives two kinds of knowledge:
an additive \emph{rule}, which can be computed from each modality separately, and a set of \emph{exceptions}, which
cannot. Section~\ref{sec:tables} defines answer tables and how we build them for each dataset.
Section~\ref{sec:rule} defines the rule and the exceptions and how we compute them exactly.
Section~\ref{sec:theory} states the properties we rely on. Section~\ref{sec:measures} defines the measures used in
the experiments, and Section~\ref{sec:flow} summarises the pipeline.

%------------------------------------------------------------------------------
\section{Answer tables}\label{sec:tables}

A two-modality task pairs an input $x_1$ from modality~1 with an input $x_2$ from modality~2 and asks a yes/no question
with answer $y\in\{0,1\}$. We assume each input belongs to a discrete \emph{index}: $a(x_1)\in\mathcal{A}$ and
$b(x_2)\in\mathcal{B}$ (for example a class, a question, or an image), and that the correct answer depends only on the
pair of indices.

\begin{definition}[Cell, answer table]
A \emph{cell} is a pair $(a,b)\in\mathcal{A}\times\mathcal{B}$. The \emph{answer table} is the map
$T:\Omega\to\{0,1\}$, where $\Omega\subseteq\mathcal{A}\times\mathcal{B}$ is the set of cells on which the question is
well-posed and evaluated. The table is \emph{complete} if $\Omega=\mathcal{A}\times\mathcal{B}$ and \emph{partial}
otherwise. Each cell carries a weight $w_{ab}\ge 0$ (its share of the evaluated samples), with
$\sum_{(a,b)\in\Omega} w_{ab}=1$.
\end{definition}

\noindent
Datasets and how each table is obtained:
\begin{itemize}[leftmargin=1.5em]
  \item \textbf{Constructed tasks.} $\mathcal{A}=\mathcal{B}=$ ten CIFAR-100 classes. A sample takes view~1 of an image
  of class $a$ and view~2 of a \emph{different} image of class $b$; its label is $T(a,b)$ for a chosen $10\times10$
  table. Complete tables; $w_{ab}=1/100$. We use 125 tables covering a wide range of structure (threshold tables,
  threshold tables with flipped cells, XOR blocks, random tables).
  \item \textbf{CLEVR.} $\mathcal{A}$ = ten yes/no questions of one question family, $\mathcal{B}$ = 1{,}000 validation
  images. $T(a,b)$ is obtained by executing question $a$'s functional program on image $b$'s scene graph. Because
  ``the red cube'' is undefined on images without exactly one red cube, we use the set-lifted semantics
  (``some red cube''), which equals CLEVR's own answer on every question's original image (60{,}835/60{,}835 yes/no
  questions). Complete tables; 103 units from 28 families.
  \item \textbf{GQA.} $\mathcal{A}$ = ten existence questions of one family (\texttt{exist}, \texttt{existC},
  \texttt{existOr}, \texttt{existOrC}), $\mathcal{B}$ = 200 validation images that contain at least one object the
  questions mention. $T(a,b)$ comes from executing the question's program on the image's scene graph with object
  identifiers removed. We keep families whose execution reproduces GQA's own answers on $\ge 99\%$ of questions, and
  drop questions about superordinate categories (``animal'', ``vehicle''). Complete tables; 24 units.
  \item \textbf{POPE.} $\mathcal{A}$ = 80 COCO object categories, $\mathcal{B}$ = the 500 POPE images; $T(a,b)=1$ iff
  object $a$ is in POPE's ground-truth list for image $b$ (this reproduces all 9{,}000 POPE labels). $\Omega$ = the
  cells POPE actually asks about; $w_{ab}\propto$ the number of times POPE asks the cell.
\end{itemize}

%------------------------------------------------------------------------------
\section{The rule and the exceptions}\label{sec:rule}

\begin{definition}[Additive rule]
An \emph{additive rule} assigns a score $u_a\in\R$ to each modality-1 index and a score $v_b\in\R$ to each modality-2
index, and answers
\begin{equation}
  f_{u,v}(a,b)=\1\!\left[\,u_a+v_b>0\,\right].
\end{equation}
\end{definition}

\noindent
Any model whose decision is a threshold on $g_1(x_1)+g_2(x_2)$ is an additive rule on every realisation of the inputs:
late fusion by summed logits, a linear head on concatenated per-modality features, and a linear probe are all of this
form. A constant bias is absorbed into $u$.

\begin{definition}[Exact rule, exceptions, additive distance]
The \emph{exact rule} is the additive rule with the largest weighted accuracy on the table,
\begin{equation}\label{eq:ceiling}
  (u^\star,v^\star)\in\argmax_{u,v}\ \sum_{(a,b)\in\Omega} w_{ab}\,\1\!\left[f_{u,v}(a,b)=T(a,b)\right].
\end{equation}
Its weighted accuracy is the \emph{additive ceiling} $C(T)$, and $d_{\mathrm{add}}(T)=1-C(T)$ is the
\emph{additive distance}. The \emph{exceptions} are the cells the exact rule gets wrong,
\begin{equation}
  E=\{(a,b)\in\Omega:\ f_{u^\star,v^\star}(a,b)\neq T(a,b)\},
\end{equation}
and the \emph{rule cells} are $\Omega\setminus E$.
\end{definition}

\paragraph{Exact computation.} Problem \eqref{eq:ceiling} is a mixed-integer linear program. With an error indicator
$z_{ab}\in\{0,1\}$ per cell, a margin of $1$ and a large constant $M$,
\begin{align}
  \min_{u,v,z}\quad & \sum_{(a,b)\in\Omega} w_{ab}\,z_{ab}\\
  \text{s.t.}\quad & u_a+v_b+M z_{ab}\ \ge\ 1 && \text{if } T(a,b)=1,\nonumber\\
  & u_a+v_b-M z_{ab}\ \le\ -1 && \text{if } T(a,b)=0,\nonumber\\
  & |u_a|,|v_b|\le U,\qquad z_{ab}\in\{0,1\}.\nonumber
\end{align}
Because $f_{u,v}$ is unchanged by positive rescaling of $(u,v)$, the unit margin loses nothing once $U$ and $M$ are large
enough ($U=50$, $M=201$). Two columns $b,b'$ with the same labels and weights on every row can share one score $v_b$
without loss, so we merge identical columns and add their weights before solving; this makes the program small
(at most $2^{|\mathcal{A}|}$ distinct columns for complete tables). We solve with HiGHS. When a time limit is reached
we report the best solution found and the solver's upper bound on $C(T)$. Constructed and GQA tables solve to
optimality; on CLEVR 62\% of units solve to optimality and the mean gap to the bound is $0.009$; for POPE the bound is
reported with every number.

\paragraph{Why not a trained additive model.} Logistic regression on one-hot indices, any gradient-trained additive
model, and EMAP all optimise a smooth surrogate (log-loss or squared error), not the 0/1 objective in
\eqref{eq:ceiling}. Their accuracy is a lower bound on $C(T)$ and was $4$--$5$ points below it on CLEVR and GQA. Only
the exact program gives the ceiling and the exceptions.

\paragraph{Graded misfit.} The exact rule gives a yes/no label per cell. For a graded score we use the marginal answer
rates under the cell weights,
\[
  p(y)=\sum_{(a,b):\,T(a,b)=y}w_{ab},\qquad p(y\mid a)=\frac{\sum_{b:\,T(a,b)=y}w_{ab}}{\sum_{b}w_{ab}},\qquad p(y\mid b)\ \text{likewise},
\]
and define
\begin{equation}\label{eq:misfit}
  m(a,b)\;=\;-\log\frac{p(y\mid a)\,p(y\mid b)}{p(y)}\,\Big|_{y=T(a,b)}
  \;=\;\log p(y)-\log p(y\mid a)-\log p(y\mid b).
\end{equation}
The fraction inside the logarithm is the (unnormalised) probability that an additive model in log-odds, built from the
row and column answer rates alone, gives to the correct answer, so $m$ is large where the additive account fits the cell
badly. For a deterministic answer, $m(a,b)$ equals the pointwise co-information
$i(x_1;x_2;y)$ with a sign flip, i.e.\ the local synergy-minus-redundancy of partial information decomposition.

%------------------------------------------------------------------------------
\section{Properties used in the analysis}\label{sec:theory}

\begin{definition}[XOR square]
Rows $a\neq a'$ and columns $b\neq b'$ form an \emph{XOR square} if $T(a,b)=T(a',b')\neq T(a,b')=T(a',b)$.
\end{definition}

\begin{theorem}[Rule without exceptions]\label{thm:ferrers}
For a complete table the following are equivalent: (i) $E=\varnothing$; (ii) $T$ contains no XOR square;
(iii) $T$ is the only $0/1$ table with its row sums and column sums.
\end{theorem}
\noindent\emph{Source.} (i)$\Leftrightarrow$(ii) is the characterisation of Ferrers relations (biorders);
(ii)$\Leftrightarrow$(iii) follows from Ryser's interchange theorem, since swapping an XOR square for its mirror image
keeps every row and column sum and every two tables with equal sums are linked by such swaps.

\begin{proposition}[Additive ceiling]\label{prop:ceiling}
Let $x_1$ and $x_2$ be drawn independently given their indices. Then every additive readout
$\1[g_1(x_1)+g_2(x_2)>0]$ has expected accuracy at most $C(T)$.
\end{proposition}
\noindent\emph{Proof sketch.} Fix one draw of $x_1$ per index $a$ and one draw of $x_2$ per index $b$ and set
$u_a=g_1(x_1)$, $v_b=g_2(x_2)$; the result is an additive rule, so its accuracy is at most $C(T)$. Averaging over draws
gives the expected accuracy of the readout. \qed

\begin{proposition}[PID cannot locate exceptions]\label{prop:pid}
For deterministic answers, every BROJA PID atom depends on $T$ only through the per-modality statistics
$p(a,y)$ and $p(b,y)$. On partial tables, two tasks with identical atoms can have $E=\varnothing$ and $E\neq\varnothing$;
on complete tables, tasks with identical atoms can have different $d_{\mathrm{add}}$.
\end{proposition}
\noindent
Theorem~\ref{thm:ferrers} and Proposition~\ref{prop:ceiling} motivate the split: the rule is fixed by one number per
row and per column, so any data that touches every row and column constrains it; an exception is constrained only by
data from that cell, and no additive component can represent it.

%------------------------------------------------------------------------------
\section{Measures}\label{sec:measures}

All measures are computed on held-out samples and split by cell type $g\in\{\text{rule},\text{exception}\}$.

\paragraph{Continual-learning protocol.} A model with one encoder per modality and a fusion head is trained on task~1
(30 epochs) and then on task~2 (10 epochs), where task~2 is another table over different classes (constructed tasks)
or another question family on the same images (CLEVR, GQA), or the COCO object questions POPE does not ask (POPE). For each method $\mathcal{M}$ we record, for every task-1 test
sample $i$, whether it is correct before task~2 ($c_i^0$) and after ($c_i^{\mathcal{M}}$). Methods: finetune, EWC,
LwF, feature distillation, experience replay (ER) and DER++, with replay buffers selected at random, by herding, by low
confidence, by forgetting events, from exception cells only, or half from exception cells and half at random
(\emph{rule + exceptions}). Regularisation strengths are chosen per unit as the setting with the best retention among
those whose task-2 accuracy is within $0.02$ of finetune (matched plasticity).

\paragraph{Accuracy, backward transfer, retention.} With $I_g$ the task-1 test samples of type $g$,
\begin{align}
  \mathrm{Acc}_g^{\,0}&=\tfrac{1}{|I_g|}\textstyle\sum_{i\in I_g} c_i^0, &
  \mathrm{Acc}_g^{\mathcal{M}}&=\tfrac{1}{|I_g|}\textstyle\sum_{i\in I_g} c_i^{\mathcal{M}}, &
  \mathrm{BWT}_g^{\mathcal{M}}&=\mathrm{Acc}_g^{\mathcal{M}}-\mathrm{Acc}_g^{\,0},\\
  \mathrm{Ret}_g^{\mathcal{M}}&=\frac{\sum_{i\in I_g} c_i^0\,c_i^{\mathcal{M}}}{\sum_{i\in I_g} c_i^0}, &
  \mathrm{ACC}^{\mathcal{M}}&=\tfrac12\big(\mathrm{Acc}^{\mathcal{M}}_{\text{task 1}}+\mathrm{Acc}^{\mathcal{M}}_{\text{task 2}}\big).
\end{align}
BWT is the standard backward transfer; retention conditions on what was known before task~2. Reporting both by cell
type matters because exceptions are a small share of samples and do not move the aggregate.

\paragraph{Exception-aware transfer metrics.} Standard BWT and FWT average over samples, so a method can score well
while losing exactly the samples that need fusion. We therefore define, for each type $g$:
\begin{align}
  \mathrm{KS}_g^{\mathcal{M}}&=\frac{\mathrm{Acc}_g^{\mathcal{M}}-\tfrac12}{\mathrm{Acc}_g^{\,0}-\tfrac12}
  &&\text{(kept share: the fraction of above-chance task-1 knowledge kept)},\\
  \mathrm{FWT}_g^{\mathcal{M}}&=\mathrm{Acc}_{2,g}^{\mathcal{M}}-\mathrm{Acc}_{2,g}^{\,\text{scratch}}
  &&\text{(task-2 accuracy on type $g$ vs.\ a model trained on task~2 alone, same budget)},\\
  \Delta\mathrm{KS}^{\mathcal{M}}&=\mathrm{KS}_{\text{exc}}^{\mathcal{M}}-\mathrm{KS}_{\text{rule}}^{\mathcal{M}},
  \qquad
  \Delta\mathrm{FWT}^{\mathcal{M}}=\mathrm{FWT}_{\text{exc}}^{\mathcal{M}}-\mathrm{FWT}_{\text{rule}}^{\mathcal{M}}, \\
  \mathrm{ACC}_{\text{bal}}^{\mathcal{M}}&=\tfrac12\textstyle\sum_{t\in\{1,2\}}\tfrac12\big(\mathrm{Acc}_{t,\text{rule}}^{\mathcal{M}}+\mathrm{Acc}_{t,\text{exc}}^{\mathcal{M}}\big)
  &&\text{(type-balanced accuracy)}.
\end{align}
$\mathrm{KS}_g$ is $1$ when nothing is lost and $0$ at chance, so it compares types that start at different accuracy;
it is left undefined when $\mathrm{Acc}_g^{\,0}<0.55$. $\Delta\mathrm{KS}<0$ means the method forgets exceptions faster
than the rule (backward), and $\Delta\mathrm{FWT}<0$ means it hinders learning new exceptions more than the new rule
(forward; a per-type form of intransigence). Task-2 cell types come from the exact rule of the task-2 answer table.

\paragraph{Matched-confidence gap.} Exceptions are often harder, so a raw difference could reflect difficulty. Let
$s$ index strata formed by unit and decile of the task-1 model's confidence $|\ell_1-\ell_0|$ (logit margin), and $n_s$
the number of exception samples correct before task~2 in stratum $s$. The gap is
\begin{equation}
  \Delta^{\mathcal{M}}=\frac{\sum_s n_s\left(\mathrm{Ret}_{\text{exc},s}^{\mathcal{M}}-\mathrm{Ret}_{\text{rule},s}^{\mathcal{M}}\right)}{\sum_s n_s},
\end{equation}
using strata with at least three samples of each type.

\paragraph{Repair after forgetting.} After task~2 the encoders are frozen and a new head is trained on 20{,}000
task-1 samples: a linear head (additive, the \emph{cheap repair}) or a two-layer MLP head (able to fuse). Among
forgotten samples ($c^0_i=1$, $c^{\mathcal{M}}_i=0$) we report the recovered share $r_g$ and its chance-corrected form
$\tilde r_g=(r_g-\tfrac12)/\tfrac12$, which is $0$ for guessing and negative when the head is systematically wrong.

\paragraph{Predicting unrepairable samples.} For each forgotten sample we score its cell by the exception indicator,
by the graded misfit \eqref{eq:misfit}, and by low confidence $-|\ell_1-\ell_0|$, and report the AUC for predicting
that the linear repair does not recover it.

\paragraph{Whether a model fuses.} For a trained model with output $h(x_1,x_2)$ on a sample set $\{(x_1^{(i)},x_2^{(i)})\}$,
the EMAP projection is
\begin{equation}
  \hat h(x_1^{(i)},x_2^{(i)})=\tfrac1n\textstyle\sum_j h(x_1^{(i)},x_2^{(j)})+\tfrac1n\sum_j h(x_1^{(j)},x_2^{(i)})-\tfrac1{n^2}\sum_{j,k}h(x_1^{(j)},x_2^{(k)}),
\end{equation}
the closest additive function in squared error. The accuracy drop from $h$ to $\hat h$ measures how much the model uses
interactions; it answers a question about the model, whereas $C(T)$ answers one about the task.

\paragraph{Statistics.} Units (one table, one CLEVR question group) are the resampling unit: all intervals are 95\%
bootstrap intervals over units. Every test, threshold and analysis script was fixed and hashed before its run.

%------------------------------------------------------------------------------
\section{Pipeline}\label{sec:flow}

\begin{enumerate}[leftmargin=1.5em, label=\textbf{Step \arabic*.}]
  \item \textbf{Build the answer table.} Choose $\mathcal{A}$, $\mathcal{B}$ and the well-posed cells $\Omega$; obtain
        $T$ from the table design, program execution on scene graphs, or ground-truth object lists; set weights
        $w_{ab}$ from how often each cell is sampled.
  \item \textbf{Solve for the exact rule.} Merge identical columns, solve the program for $(u^\star,v^\star)$, and
        record $C(T)$, its upper bound, and the exception set $E$. Compute the graded misfit $m(a,b)$.
  \item \textbf{Label samples.} Give every sample the type of its cell (rule or exception) and its misfit; labels on
        training cells are used only to select buffers, labels on test cells only for evaluation.
  \item \textbf{Train and evaluate.} Run task~1, then task~2 under each method; record correctness before and after
        for every task-1 test sample, and task-2 accuracy.
  \item \textbf{Measure by cell type.} Report accuracy, BWT and retention for rule and exception samples, the
        matched-confidence gap, repair recovery, prediction AUCs, and EMAP where a model's fusion is in question.
\end{enumerate}

\end{document}
